\scanpage{45}
3.1.\deleted{\unsure{13}}15. On peut \deleted{\unsure{encore}} décrire $Rf^!$ par le procédé de 3.1.13. \deleted{\unsure{Résultat}} Soit en effet un morphisme compactifiable $f$ :
\diagram{\begin{tikzcd}[ampersand replacement=\&]
 X \arrow[r,hook,"j"'] \arrow[d,"f"'] \& \bar X \arrow[dl,"\bar f"]\\
 S \& {}
\end{tikzcd}}
et $\mathcal A$ un faisceau d'anneaux sur $S$ tel que $n\mathcal A=0$ $(n\geq1)$. Pour $W=(U,V,\unsure{\varphi|\psi})\in\mathrm{Ob}\,X_f$ (notation de 3.1.13), soit
\[
 K(W)=\bar f_{V*}\tau_{\leq 2d}C^*_{\ell,\bar X_V}(\mathcal A_U)
 \tag*{(3.1.15.1)}
\]
\margin{$\tau''$ ?}
% slanted
\margin{signification pas claire}
(résolution calculée sur $\bar X_V=\bar X\times_S V$). \deleted{On a un quasi-isomorphisme}
\[
 \deleted{j_!K(W)}
\]
% diagram in the left margin, with the diagonal arrows directed down and left
\diagram{\begin{tikzcd}[ampersand replacement=\&,row sep=small,column sep=small]
 \bar X \arrow[d] \arrow[dd,bend right=48,"\bar f"'] \& \bar X_V \arrow[d] \arrow[dd,bend right=42,"\bar f_V"'] \& {}\\
 X \arrow[d,"f"'] \& X_V \arrow[l] \arrow[d,"f_V"] \& U \arrow[l]\\
 S \& V \arrow[l,"j"'] \& {}
\end{tikzcd}}
Pour $j$ \deleted{inclusion de}\inserted{morphisme structurel} $V$ dans $S$, on a un quasi-isomorphisme
\[
 j_!K(W)\xrightarrow{\sim}\bar f_*\tau_{\leq 2d}C^*_{\ell,\bar X}(\mathcal A_U)
\]
\margin{$(j\bar f_V)_*$ ?}
($\overline{\underline{\mathrm{XVII}}}$ 5.2.\tikz[baseline]{\draw (0,-.07) -- (.10,.40); \draw (.02,.02) -- (.10,-.06);})
\margin{réf}
D'après 3.1.9,\margin{?} le système (3.1.13)\margin{?} des $K(W)$ permet donc de calculer $Rf^!$.
% slanted
\margin{pas clair}
\deleted{Ceci donne un sens à la ``\uline{formule d'induction}''}
% archive sheet number: 44
